\underline{Exercise $2.4$} 

\begin{enumerate}
\item[5.] We are given the function $y=5+3x$ with domain $X = \{x | 1 \leq x \leq 9 \}$. Note that for this function, when $x=1$, $y=8$ and when $x=9$, $y=32$. So the range for this function is:
\[
f(X) =\{y \mid 8 \leq y \leq 32\}
\]
Note: It is not always the case that extreme values of the domain correspond to extreme values of the range. For example, consider $y=x^{2}$ with domain $\{x \mid-2 \leq x \leq 2\}$, the range here is $\{y \mid 0 \leq y \leq 4\}$.

\item[7.] (a) No, (b) Yes
\item[8.] For each output, we would want to produce at the lowest cost. \\
\end{enumerate}

%%%%%%%%%%%%%%%%%%%%%%%%%%%%%%%%% Exercise 2.5
\newpage

\underline{Exercise $2.5$}

% Shared look for the three line plots: axes through the origin, the axis
% labels parked at the far end of each axis (ticklabel* cs needs compat>=1.3;
% these are the only pgfplots pictures in this file, so setting it here is
% safe), and the intercepts marked.
\pgfplotsset{compat=1.18}
\definecolor{plotline}{HTML}{912040}
\pgfplotsset{
  hwline/.style={
    axis lines=middle,
    width=6.6cm, height=3.6cm,
    scale only axis,
    xlabel={$x$}, ylabel={$y$},
    every axis x label/.style={at={(ticklabel* cs:1.0)}, anchor=west},
    every axis y label/.style={at={(ticklabel* cs:1.0)}, anchor=south},
    label style={font=\small}, tick label style={font=\small},
    axis line style={-latex, line width=0.3pt},
    tick style={line width=0.3pt},
    enlargelimits=false,
  },
  hwplot/.style={color=plotline, line width=0.5mm, samples=2},
  hwdot/.style={only marks, mark=*, mark size=1.8pt, color=plotline},
}

\begin{enumerate}
\item Graph the following functions and find their inverse.
\begin{enumerate}\renewcommand{\labelenumi}{\alph{enumi})}

\item $y = 16 + 2x$, \quad $f^{-1}(y) = \dfrac{y-16}{2}$

\nopagebreak
\tagstructbegin{tag=Figure,alt={Graph of the straight line y equals 16 plus 2x: an upward-sloping line with slope 2, crossing the y-axis at 16 and the x-axis at negative 8}}\tagmcbegin{}
\begin{tikzpicture}
\begin{axis}[hwline, xmin=-11, xmax=3, ymin=-6, ymax=22,
             xtick={-8,-4}, ytick={8,16}]
\addplot [hwplot, domain=-11:3] {16 + 2*x};
\addplot [hwdot] coordinates {(-8,0) (0,16)};
\end{axis}
\end{tikzpicture}
\tagmcend\tagstructend

\item $y = 8 - 2x$, \quad $f^{-1}(y) = \dfrac{8-y}{2}$

\nopagebreak
\tagstructbegin{tag=Figure,alt={Graph of the straight line y equals 8 minus 2x: a downward-sloping line with slope negative 2, crossing the y-axis at 8 and the x-axis at 4}}\tagmcbegin{}
\begin{tikzpicture}
\begin{axis}[hwline, xmin=-2, xmax=7, ymin=-6, ymax=12,
             xtick={2,4,6}, ytick={4,8}]
\addplot [hwplot, domain=-2:7] {8 - 2*x};
\addplot [hwdot] coordinates {(4,0) (0,8)};
\end{axis}
\end{tikzpicture}
\tagmcend\tagstructend

\item $y = 2x + 12$, \quad $f^{-1}(y) = \dfrac{y-12}{2}$

\nopagebreak
\tagstructbegin{tag=Figure,alt={Graph of the straight line y equals 2x plus 12: an upward-sloping line with slope 2, crossing the y-axis at 12 and the x-axis at negative 6}}\tagmcbegin{}
\begin{tikzpicture}
\begin{axis}[hwline, xmin=-9, xmax=3, ymin=-6, ymax=18,
             xtick={-6,-3}, ytick={6,12}]
\addplot [hwplot, domain=-9:3] {2*x + 12};
\addplot [hwdot] coordinates {(-6,0) (0,12)};
\end{axis}
\end{tikzpicture}
\tagmcend\tagstructend

\end{enumerate}
\end{enumerate}

%%%%%%%%%%%%%%%%%%%%%%%%%%%%%%%%% Exercise 4.2


\underline{Exercise 4.2}  
% Question 6
\begin{enumerate}
\item[6.]
\begin{multicols}{2}
\begin{enumerate}\renewcommand{\labelenumi}{\alph{enumi})}
\item $x_{2}+x_{3}+x_{4}+x_{5}$
\item $a_{5} x_{5}+a_{6} x_{6}+a_{7} x_{7}+a_{8} x_{8}$
\item $b x_{1}+b x_{2}+b x_{3}+b x_{4}$
\item $a_{1}+a_{2} x+a_{3} x^{2}+\ldots+a_{n} x^{n-1}$
\item $x^{2}+(x+1)^{2}+(x+2)^{2}+(x+3)^{2}$  

\vspace{1em}
\end{enumerate}
\end{multicols}

% Question 8
\item[8.] 
\begin{enumerate}
\item \[ \left(\sum_{i=0}^{n} x_{i} \right)+x_{n+1} =x_{0}+x_{1}+x_{2}+\ldots+x_{n+1}  =\sum_{i=0}^{n+1} x_{i} \]
\item \[ \begin{aligned} 
\sum_{j=1}^{n} a b_{j} y_{j}&=a b_{1} y_{1}+a b_{2} y_{2}+\ldots+a b_{n} y_{n} \\ 
&= a\left(b_{1} y_{1}+b_{2} y_{2}+\ldots+b_{n} y_{n}\right) \\
&= a \sum_{j=1}^{n} b_{j} y_{j} \end{aligned}\]
\item \[ \begin{aligned}  \sum_{j=1}^{n}\left(x_{j}+y_{j}\right) 
& =\left(x_{1}+y_{1}\right)+\left(x_{2}+y_{2}\right)+\ldots+\left(x_{n}+y_{n}\right) \\
& =x_{1}+x_{2}+\ldots+x_{n}+y_{1}+y_{2}+\ldots+y_{n} \\
& =\sum_{j=1}^{n} x_{j}+\sum_{j=1}^{n} y_{j} \end{aligned} \]
\end{enumerate}
\end{enumerate}

%%%%%%%%%% Exercise 5.1


\underline{Exercise 5.1}  
\begin{enumerate}
\item[1.] \begin{multicols}{3}
\begin{enumerate}\renewcommand{\labelenumi}{\alph{enumi})}
\item $ q \implies p $
\item $ q \implies p $
\item $ q \iff p $
\item $ q \iff p $
\item $ q \iff p $
\item $ p \implies q $
\item $ q \implies p $
\end{enumerate}
\end{multicols}
\end{enumerate}
